% 第四章结果与讨论草稿：按“基线与无示例→B1示例学习→B2端到端选择→案例与讨论”重写。
% 无示例比较均按冻结的 60 例正式评估集统计。
\section{结果与讨论}

本章在冻结的 60 个事件组合上比较 B0、B1 和 B2。讨论的对象不是具体设备故障是否发生，而是系统能否获得与人工检查意图一致的候选证据，并据此形成一项可复核的固定检查建议。章节依次考察无示例条件下证据覆盖与行动选择之间的差距、固定示例对 B1 的作用、全库 LLM 证据选择的端到端表现，以及代表性案例和实施边界。

\subsection{基线与无示例条件：证据覆盖不等于正确行动}

本节比较未向 RAG-LLM 提供固定示例上下文时，各基线工作流的证据选择与核查建议生成结果。B0 使用 BM25 Top-K 检索与标签机械投票；B1 使用与 B0 完全相同的候选证据包，仅将投票替换为 LLM 行动选择；B2 则要求 LLM 从完整的 300 张人工证据卡中选择候选证据，再输出行动建议。三种配置分别对应词汇检索加规则映射、固定检索条件下的零样本行动推理，以及零样本端到端证据选择。

表~\ref{tab:zero-shot-main} 将候选证据语义质量、语义--行动对质量和最终行动表现并列报告。该比较的重点不在于将所有配置压缩为单一分数，而在于检验候选证据覆盖能否真正转化为正确的检查建议。B0 中，BM25 候选包在 \(K\geq5\) 时达到完整语义候选覆盖，但标签投票的 Accuracy 仍仅为 0.2500--0.2667。B1 零样本使用完全相同的候选证据包；随着 \(K\) 增大，其行动 Accuracy 从 0.2333 提升至 0.4833，但仍未稳定地区分四类检查建议。B2 零样本在各 \(K\) 下均达到完整候选覆盖，语义 F1 也高于 B0，却仍维持较低的 Pair-F1 和行动表现，表明全库访问本身不等同于有效的证据--行动映射。

\begin{table*}[t]
\centering
\caption{无固定示例上下文时 B0、B1 与 B2 在冻结 60 例评估集上的比较。}
\label{tab:zero-shot-main}
\small
\begin{tabular}{c c c c c c c}
\hline
配置 & \(K\) & 语义 F1@K & 候选覆盖率 & Pair-F1@K & Accuracy & Macro-F1 \\
\hline
B0 & 3  & 0.4188 & 0.9000 & 0.1698 & 0.2500 & 0.1667 \\
B0 & 5  & 0.4749 & 1.0000 & 0.1458 & 0.2500 & 0.1744 \\
B0 & 8  & 0.5325 & 1.0000 & 0.1692 & 0.2667 & 0.2079 \\
B0 & 10 & 0.5374 & 1.0000 & 0.1665 & 0.2500 & 0.1974 \\
\hline
B1 零样本 & 3  & 0.4188 & 0.9000 & 0.1587 & 0.2333 & 0.2258 \\
B1 零样本 & 5  & 0.4749 & 1.0000 & 0.1860 & 0.3167 & 0.2551 \\
B1 零样本 & 8  & 0.5325 & 1.0000 & 0.2701 & 0.4333 & 0.3401 \\
B1 零样本 & 10 & 0.5374 & 1.0000 & 0.3063 & 0.4833 & 0.3686 \\
\hline
B2 零样本 & 3 & 0.5180 & 1.0000 & 0.1708 & 0.2667 & 0.1652 \\
B2 零样本 & 5 & 0.5875 & 1.0000 & 0.1765 & 0.2500 & 0.1190 \\
B2 零样本 & 8 & 0.5719 & 1.0000 & 0.1719 & 0.2500 & 0.1210 \\
\hline
\end{tabular}
\end{table*}

图~\ref{fig:zero-shot-evidence-action} 建议以语义 F1@K、语义--行动对 F1@K 与 Accuracy 为三组并列柱状图。若候选覆盖率或语义 F1 已经较高，但 Pair-F1 和 Accuracy 仍然较低，则可直接说明：候选证据存在并不足以保证系统能够作出正确行动选择。这一现象构成后续引入固定示例的动机。

\begin{figure}[t]
  \centering
  \includegraphics[width=\columnwidth]{fig_zero_shot_evidence_action}
  \caption{B0、B1 零样本与 B2 零样本的证据质量和行动质量比较。候选证据覆盖不等于证据--行动映射正确。}
  \label{fig:zero-shot-evidence-action}
\end{figure}

\subsection{B1：固定候选证据条件下的示例学习与收益平台}

在 B1 中，候选证据始终由同一 BM25 检索器给出，因此固定示例带来的变化主要反映行动选择阶段的上下文学习。图~\ref{fig:b1-k-n} 展示不同 \(K\) 和示例数 \(n\) 下的 Accuracy 与 Macro-F1；每个点使用三个冻结示例池的均值和标准差。固定示例加入后，B1 的行动表现相较零样本显著提升，表明模型能够从完整案例中学习事件组合、候选证据与检查建议之间的任务关系。

表~\ref{tab:b1-main-settings} 给出代表性设置。K=5、\(n=8\) 的三池平均 Accuracy 为 \(0.9389\pm0.0255\)，而在 \(K=8,n=16\) 时，Accuracy 和 Macro-F1 分别达到 \(0.9889\pm0.0192\) 与 \(0.9888\pm0.0193\)。当示例数进一步增加至 \(n=20\) 时，两项指标保持不变，但平均总 token/池从约 2.739M 增至约 3.364M。该结果表明，固定示例并非越多越好：在当前任务中，性能随 \(n\) 增长而改善，但会在一定规模后进入收益平台。

\begin{figure*}[t]
  \centering
  \includegraphics[width=0.86\textwidth]{fig_b1_k_n_tradeoff}
  \caption{B1 固定示例实验中，\(K\) 与 \(n\) 对 Accuracy、Macro-F1 和总 token 的影响。误差线表示三个冻结示例池的标准差。}
  \label{fig:b1-k-n}
\end{figure*}

\begin{table}[t]
\centering
\caption{B1 代表性设置的跨池结果。}
\label{tab:b1-main-settings}
\small
\begin{tabular}{c c c c c}
\hline
\(K\) & \(n\) & Accuracy & Macro-F1 & 平均总 token/池 \\
\hline
5 & 8  & \(0.9389\pm0.0255\) & \(0.9387\pm0.0256\) & [待填] \\
8 & 16 & \(0.9889\pm0.0192\) & \(0.9888\pm0.0193\) & 约 2.739M \\
8 & 20 & \(0.9889\pm0.0192\) & \(0.9888\pm0.0193\) & 约 3.364M \\
\hline
\end{tabular}
\end{table}

这一结果不意味着模型已形成故障诊断能力，而是说明在固定、可复现的候选证据边界内，示例学习能够更稳定地将事件组合映射为一项固定人工检查建议。B1 因而构成本文的主结果路径：BM25 提供稳定的候选依据，LLM 在受限行动空间中完成证据条件下的优先核查选择。

\subsection{B2：全库 LLM 证据选择的端到端消融}

B2 进一步要求 LLM 面向完整证据库选择候选证据。与 B1 相比，该配置不再将候选证据质量固定在检索阶段，因此同时引入证据选择误差和行动选择误差。零样本结果已表明全库访问本身不会自动提高表现；本文随后通过 \textit{selector}、\textit{action} 和 \textit{both} 三种设置，分析固定示例应作用于证据选择端、行动选择端还是两个阶段。

表~\ref{tab:b2-stage-ablation} 以 \(K=8,n=8\) 的 pool 01 为例。仅在证据选择端加入示例时，Accuracy 为 0.3167；仅在行动选择端加入示例时，Accuracy 提升至 0.7500；两个阶段同时加入示例时，Accuracy 为 0.9167，Macro-F1 为 0.9169。该结果表明，行动选择端的示例支持对最终输出尤为关键，而两个阶段同时获得示例支持时，证据语义与行动标签之间的一致性进一步改善。

\begin{table}[t]
\centering
\caption{B2 阶段消融：\(K=8,n=8\)，pool 01。}
\label{tab:b2-stage-ablation}
\small
\begin{tabular}{c c c c}
\hline
配置 & Accuracy & Macro-F1 & 合同合规率 \\
\hline
selector & 0.3167 & 0.2422 & 1.0000 \\
action   & 0.7500 & 0.7442 & 1.0000 \\
both     & 0.9167 & 0.9169 & 1.0000 \\
\hline
\end{tabular}
\end{table}

然而，B2 的改善并未使其整体超过 B1 的固定候选证据主设置。在 \(K=8,n=8\) 的三个冻结示例池中，B2-\textit{both} 的 Accuracy 分别为 0.9167、0.8333 和 0.8833，平均为 \(0.8778\pm0.0419\)；Macro-F1 平均为 \(0.8777\pm0.0413\)。这说明 LLM 可在端到端配置中学习部分证据选择规律，但完整证据库中的自主筛选仍带来额外误差与示例池波动。图~\ref{fig:b2-stage-metrics} 应同时呈现语义证据 F1、语义--行动对 F1 与 Accuracy，以避免将最终行动分数误解为全部证据选择质量。

\begin{figure}[t]
  \centering
  \includegraphics[width=\columnwidth]{fig_b2_stage_metrics}
  \caption{B2 阶段消融与跨池结果中的证据语义质量、语义--行动对 F1 和最终行动 Accuracy。}
  \label{fig:b2-stage-metrics}
\end{figure}

本文未将 B2 固定示例网格继续扩展至 \(K=10\)。在 B2 零样本的 \(K=10\) 运行中，行动输出合同合规率降至 0.4333，并出现 41 个行动阶段错误。这表明完整证据库选择与更大候选包共同进入行动上下文时，输出稳定性会明显下降。因此，B2 的固定示例消融限定在 \(K=3,5,8\)，以考察可稳定执行的端到端配置，而不是无约束地扩大上下文。

\subsection{代表性案例与结果讨论}

平均指标揭示总体趋势，但不能显示错误发生在证据获取还是行动选择阶段。表~\ref{tab:case-analysis} 建议选取四类代表性案例：B0 中候选证据覆盖检查意图但标签投票错误的案例；B1 在相同证据包上经示例学习后纠正行动的案例；B2 选择到部分相关证据但行动仍错误的案例；以及 B2-\textit{both} 同时获得可接受证据与正确行动的案例。表中仅报告事件组合摘要、候选证据的语义类别、模型行动与金标准行动，不将异常现象解释为具体故障原因。

\begin{table*}[t]
\centering
\caption{代表性案例的证据--行动链分析。正式写作时从保存的个案 JSON/JSONL 输出中填入。}
\label{tab:case-analysis}
\small
\begin{tabular}{p{0.15\textwidth} p{0.25\textwidth} p{0.24\textwidth} p{0.16\textwidth} p{0.16\textwidth}}
\hline
案例类型 & 事件组合摘要 & 候选证据语义 & 模型行动 & 金标准行动 \\
\hline
B0 标签投票错误 & [待填] & [覆盖正确意图] & [待填] & [待填] \\
B1 示例后纠正 & [待填] & [与 B0 相同] & [待填] & [待填] \\
B2 部分证据命中但行动错误 & [待填] & [部分覆盖] & [待填] & [待填] \\
B2-\textit{both} 成功案例 & [待填] & [可接受证据语义] & [待填] & [待填] \\
\hline
\end{tabular}
\end{table*}

综合而言，B0、B1 和 B2 的分层比较确认：在中子电流异常预检任务中，候选证据的获取方式与证据条件下的行动选择同样重要。B0 与 B1 使用完全相同的 BM25 候选证据包，而 B1 通过 LLM 替代标签机械投票后取得更高的行动表现，说明词汇检索获得的证据仍需要与事件组合进行受约束的证据--行动推理。B2 进一步表明，全库证据选择虽然增加了端到端任务难度，但在证据选择端和行动选择端同时加入固定示例后，系统能够显著改善证据--行动一致性。

与仅使用 BM25 标签投票的 B0、以及由 LLM 面向全库自主选择证据的 B2 配置相比，B1 的结果确认：即便语言模型能够访问相关技术证据，若缺少稳定的候选证据获取机制或受约束的证据--行动映射，其最终检查建议仍可能表现较弱，并出现证据语义与行动标签不一致的情况。这与检索增强决策支持研究中的发现一致：领域证据只有在与具体任务表示、输出约束和人工工作流相结合时，才能稳定支持模型生成。本文框架通过在统一的事件--证据--行动结构中关联中子电流异常、可定位人工证据和固定检查建议，扩展了这一思路；它不将证据解释为故障确认，而是将其组织为支持人工优先核查的可审查依据。

上述结果均限于当前 300 张人工证据卡、冻结 60 例评估集和 \texttt{qwen3.8-flash} 模型设置。它们不构成对所有检索器或所有语言模型的普遍排名，也不构成具体设备故障的自动诊断结论。
